\section{Open Science}
\label{sec:open-science}

The artifacts behind the formal and empirical claims are released for
public inspection. This appendix names what is available, how to
rebuild it, and how a reviewer recovers each load-bearing statement
from its primary source.

\subsection{Artifact availability}
\label{sec:open-science:availability}

Three artifact classes accompany the paper. The Lean~4 mechanization
that discharges the four theorems of Section~\ref{sec:model} is the
central artifact: a single paper-local source module of 589 lines,
kept beside the paper rather than in the substrate's proof root, and
packaged as a self-contained Lake project together with the three
substrate Treaty modules that form its import closure. The packaged
copy is byte-identical to the module of record, and the supplementary
package records the SHA-256 that a reviewer checks this against. The
bibliography file lists primary-source
citations to specifications, standards documents, and peer-reviewed
papers. The paper text and the supplementary package are released under
the same permissive license that governs the substrate's public
artifacts. For double-blind review the artifacts are submitted in
anonymized form; a public repository location is named in the
camera-ready version.

\subsection{Reproducing the Lean mechanization}
\label{sec:open-science:lean}

The Lake project pins a Lean toolchain through a checked-in
\path{lean-toolchain} file, matching the release the substrate's own
proof root pins, and declares no external dependency. A reviewer checks
the mechanization from a fresh extract:

\begin{verbatim}
tar xzf lean-source.tar.gz
cd chio-lean
lake build
lake env lean Chio/Treaty/SensorGroundedAdmission.lean
\end{verbatim}

The \path{lake build} step builds the three substrate Treaty modules;
the \path{lake env lean} step elaborates the paper-local module against
them and exits with no output. The two steps stay separate because the
module is checked exactly as it sits in the repository: from outside
the root module, never registered in it. Against the substrate's own
proof root, the same check is \path{lake env lean} on
\path{lean/SensorGroundedAdmission.lean} run from
\path{formal/lean4/Chio}, and that is the run of record. Neither route
reports a warning, and the sources contain no \texttt{sorry}. The
\texttt{\#print axioms} command run against each of the four theorems
reports only the standard Lean kernel axioms \texttt{propext},
\texttt{Classical.choice}, and \texttt{Quot.sound}. No project-local
\texttt{axiom} declaration appears anywhere in the dependency tree.
The supplementary package's \path{proof-manifest.toml} records the
fully qualified theorem names, their source modules, and their axiom
sets for tool consumption, and records that the module is paper-local,
outside the substrate's own theorem inventory.

\subsection{Reproducing the empirical claims}
\label{sec:open-science:empirical}

The empirical claims of Section~\ref{sec:evaluation} rest on the
sensor-state attestation field carried by every signed receipt produced
by the admission kernel described in Section~\ref{sec:implementation}.
A reviewer reproduces the figures by exercising the kernel against the
receipt corpus described in Section~\ref{sec:implementation},
extracting each receipt's signed attestation block, and verifying both
the attestation signature and the DSSE subject digest that binds
attestation bytes to body bytes. The canonical-JSON parser and the
provider-record schema are part of the released runtime package. The
runtime's deployment identity is anonymized for review and named in
the camera-ready version.

\subsection{Bibliography accessibility}
\label{sec:open-science:bib}

Every load-bearing citation resolves to a primary source that is
publicly retrievable. The bibliography includes vendor specifications
(Intel TDX, AMD SEV-SNP, Apple Private Cloud Compute, Arm CCA, AWS
Nitro), Internet standards (RFC~8785, RFC~9334), and peer-reviewed
publications from USENIX, IEEE, and ACM venues. No paywalled
proceeding or unavailable technical report carries a load-bearing
claim.
